\section{Related work}\label{sec:related}

\subsection{Graph collaborative filtering and graph debiasing}\label{sec:related-graph}

LightGCN \citep{he2020lightgcn} is the backbone used in our experiments. Other graph
convolutional recommenders include NGCF \citep{wang2019ngcf} and UltraGCN
\citep{mao2021ultragcn}, and contrastive graph variants include SGL \citep{wu2021sgl} and
SimGCL \citep{yu2022simgcl}. Linear and closed-form models include EASE \citep{steck2019ease},
GF-CF \citep{shen2021gfcf} and BSPM \citep{choi2023bspm}. Weighted matrix factorisation for
implicit feedback, iALS \citep{hu2008ials}, is a standard baseline, and \citet{rendle2022ials}
revisit its tuning on item recommendation benchmarks. \citet{wu2022gnnsurvey} survey graph
neural recommenders. These models fit or propagate over the logged graph and do not model
which pairs were exposed.

Graph models have also been debiased. NAVIP \citep{kim2022navip} weights neighbour
aggregation by inverse propensities. r-AdjNorm \citep{zhao2022radjnorm} tunes the strength of
adjacency normalisation to trade accuracy against novelty, and \citet{chen2024gcnpop} study
how graph convolutions can amplify popularity bias. APDA \citep{zhou2023apda} and CAGED
\citep{que2025caged} learn aggregation weights, and DPAA \citep{dpaa2026}, an arXiv preprint,
re-weights message passing. AdvDrop \citep{zhang2024advdrop} learns an adversarial edge dropout
on the interaction graph and reports IPS and DR baselines on Coat, Yahoo! R3 and KuaiRec. PPD
\citep{islam2026ppd} removes a popularity direction from pre-trained LightGCN embeddings
without retraining, and is evaluated on KuaiRec, Coat and Yahoo! R3. DR-GNN
\citep{wang2024drgnn} applies distributionally robust optimisation over graph neighbourhoods;
its ``DR'' stands for distributionally robust, not doubly robust. The popularity-debiasing
methods MACR \citep{wei2021macr} and PDA \citep{zhang2021pda} are used as baselines in our
study. Most of these methods act on aggregation, adjacency, embeddings or a robust objective,
rather than on the negatives of the pairwise loss.

\subsection{Propensity, doubly robust and negative-weighting methods}\label{sec:related-dr}

Ratings and implicit feedback are missing not at random \citep{marlin2009collaborative}, and
surveys organise the resulting biases and the causal methods that address them
\citep{chen2023bias,gao2024causalsurvey}. Propensity methods re-weight the loss by inverse
propensities or self-normalised estimators \citep{schnabel2016,swaminathan2015snips,saito2020unbiased}.
\citet{steck2010mnar} showed that missing ratings carry information, proposed test measures
that can be estimated without bias when ratings are missing not at random, and derived training
objectives that account for the missing entries. Doubly robust learning adds an imputation model
to the inverse-propensity estimator \citep{wang2019drjl}, and later doubly robust variants
address related debiasing problems \citep{guo2021mrdr,dai2022gdr,li2023stabledr}. Information-theoretic counterfactual
learning from missing-not-at-random feedback is studied by \citet{wang2020cvib}.

Some debiasing methods use a small uniform sample for training. AutoDebias
\citep{chen2021autodebias} learns sample weights and an imputation function by bi-level
meta-learning on a small uniform sample. CausE \citep{bonner2018cause} embeds a small randomised
sample alongside the logged one, and KDCRec \citep{liu2020kdcrec} distils from a model fitted on
uniform data into a model trained on the biased log.

The methods closest to our negative weighting are those for exposure and false negatives. ExpoMF
\citep{liang2016exposure} infers an exposure posterior for unclicked items inside a matrix
factorisation model. Rel-MF \citep{saito2020unbiased} re-weights observations by propensity, and
unbiased pairwise learning \citep{saito2020ictir} applies the same idea to pairwise objectives.
SRNS \citep{ding2020srns} identifies likely false negatives during sampling, and debiased
contrastive learning \citep{chuang2020dcl} corrects negatives that are in fact positives. In
contrastive collaborative filtering, the BC loss \citep{zhang2022bc} adds a margin that depends
on the popularity-bias degree of each interaction, and AdvInfoNCE \citep{zhang2023advinfonce}
assigns each negative an adversarially learned hardness weight. Correct-and-Weight
\citep{yin2026cw}, an arXiv preprint, weights negatives against false negatives and is tested
with graph backbones on random splits. Exposure-aware negative sampling for graph recommenders
uses in-community popularity as a proxy for exposure \citep{chen2026icpns}, also an arXiv
preprint. The original formulations of ExpoMF, SRNS, debiased contrastive learning and
AdvInfoNCE do not combine a propensity with an imputation of the kind a doubly robust estimator
provides. IPS-weighted BPR on LightGCN has been studied with simulated propensities
\citep{raja2025ipsbpr}.

The offset used by our candidate models is a fixed, standardised imputation logit added to the
score. Adding a fixed offset is a familiar device. In long-tail classification, logit adjustment
adds a class prior to the logits \citep{menon2021logit}. The offset used here differs in one
respect: it is a debiased estimate of item-level preference, obtained from an
inverse-propensity-weighted imputation, and not a class prior.

\subsection{Evaluation and reproducibility}\label{sec:related-eval}

Several studies have shown that comparisons between recommenders depend on the protocol.
\citet{dacrema2019progress} re-ran recent neural recommenders against simple baselines and found
that several did not outperform well-tuned versions of those baselines. \citet{rendle2019baselines}
argue that baselines in published comparisons are often under-tuned, so that a reported gain
depends on the effort spent on each method. \citet{rendle2022ials} show that iALS, when tuned
carefully, is competitive on common item recommendation benchmarks. \citet{anelli2023myth}
question the reported advances of graph collaborative filtering.

Several parts of the protocol can change a conclusion. Sampled and exact metrics can disagree
\citep{krichene2020sampled}. \citet{nadeau2003} propose a corrected resampled t-test for the generalisation error, for
resamples that overlap. Reporting the expected
validation performance as a function of the number of configurations tried makes the tuning
budget visible \citep{dodge2019showyourwork}. Comparisons over several methods and datasets also
need a multiple-testing correction, for which we use the step-down procedure of
\citet{holm1979}.

\subsection{Gap and positioning}\label{sec:related-gap}

The contribution of this paper is evidence rather than a new method. The graph debiasers and the
candidate models of Section~\ref{sec:methods} are evaluated under one protocol, and the study measures how far
each conclusion depends on evaluation choices. Table~\ref{tab:positioning} summarises how the
related studies use uniformly collected data and which benchmarks they use. The uses differ.
Some studies train on uniform data (AutoDebias, CausE, ConAlign), some select configurations on
part of it (CDR, PPD, ConAlign), and some keep it out of both training and selection (DCE, BISER).

\begin{table}[t]
\centering
\caption{Positioning of related studies by their protocol. Entries follow the protocols as
described in the cited papers; ``not stated'' marks a point that the description does not
settle. Uniform data are the uniformly collected or randomised samples of the benchmark.}
\label{tab:positioning}
\footnotesize
\adjustbox{max width=\textwidth}{%
\begin{tabular}{p{2.6cm}p{1.9cm}p{2.5cm}p{2.0cm}p{2.6cm}p{2.4cm}}
\toprule
Study & Graph backbone & Propensity or imputation & Uniform data in training & Uniform data for validation or selection & Unbiased test benchmarks \\
\midrule
AutoDebias \citep{chen2021autodebias} & no (matrix factorisation) & learned weights and imputation & yes, small sample & yes, small sample & Coat, Yahoo! R3 \\
CausE \citep{bonner2018cause} & no (matrix factorisation) & none & yes, randomised sample & no, validation from logged policy & MovieLens, Netflix \\
CDR \citep{song2023cdr} & no (matrix factorisation) & doubly robust family & no & yes, part of unbiased data & Coat, Yahoo! R3, KuaiRand \\
DCE \citep{kweon2024dce} & no (MF or NCF) & doubly robust loss & no & no, held out from biased training data & Coat, Yahoo! R3, KuaiRec \\
BISER \citep{lee2022biser} & no & self-inverse propensity & no & no, held out from biased training data & Coat, Yahoo! R3 \\
ConAlign \citep{zhang2026conalign} & no (dual tower) & MF-IPS and MF-DR baselines & yes & yes & Coat, Yahoo! R3, KuaiRand-Pure \\
NAVIP \citep{kim2022navip} & yes & inverse propensity on edges & no & pseudo-unbiased popularity sample & Yahoo! R3 and pseudo-unbiased sets \\
AdvDrop \citep{zhang2024advdrop} & yes & IPS and DR baselines only & no, for its own method & not stated & Coat, Yahoo!, KuaiRec \\
PPD \citep{islam2026ppd} & yes, post hoc & none & no & yes on Coat and KuaiRec; biased sample on Yahoo! R3 & Coat, Yahoo! R3, KuaiRec \\
DPAA \citep{dpaa2026} & yes & none & no & validation present, construction not confirmed & Coat, Yahoo! R3, KuaiRec \\
DR-GNN \citep{wang2024drgnn} & yes & none & not stated & not stated & Coat, Yahoo \\
\midrule
This study & yes & IPS-fitted imputation and propensity & no, except Coat shipped propensities & yes & Coat, Yahoo! R3, KuaiRand-Pure, KuaiRec \\
\bottomrule
\end{tabular}}
\end{table}

A web search in October 2026 found no graph recommender that uses a posterior probability,
built from a propensity and an imputation, to weight sampled negatives, and none that adds a
debiased imputation to a graph model's score. It also found no graph study that uses
KuaiRand-Pure. The search used general web sources and cannot be exhaustive. Coat, Yahoo! R3 and
KuaiRec are already used for graph debiasing by AdvDrop, PPD and DPAA. We therefore do not claim
that the posterior weighting or the imputation offset is new. The claim we make is narrower: we
found no study that evaluates graph debiasers on all four benchmarks under one protocol, with
uniform data excluded from training, and with an analysis of how the budget and the selection
rule change the conclusions.
